\chapter{分片、复制与共识}

分布式系统通过分片扩展容量，通过复制提高可用性，通过共识在故障中维护关键状态一致。三者共同构成大多数存储和计算平台的基础。

\section{分片}

分片把数据按 key、范围或哈希分到不同节点。好的分片要负载均衡、支持扩容、减少跨分片事务。热点 key 会破坏均衡，需要拆分、缓存、限流或特殊路由。

分片后的查询可能需要 scatter-gather：向多个分片发请求，再合并结果。它类似并行归约，但网络延迟和节点故障让合并更复杂。

分片策略要匹配访问模式。哈希分片适合点查和均匀 key，范围分片适合范围扫描和排序，但容易出现热点区间。虚拟分片把逻辑分片数量做得比物理节点多，便于迁移和均衡。热点 key 不能只靠增加节点解决，因为所有请求仍然打到同一个逻辑位置；需要缓存、拆 key、限流、异步化或改变业务模型。

\section{再分片}

系统运行一段时间后，数据量和热点会变化，需要再分片。再分片要迁移数据、更新路由、处理迁移期间的读写，并保证失败后可恢复。简单哈希分片在节点数变化时可能移动大量 key，一致性哈希或虚拟分片可以减少迁移范围。

再分片最难的是迁移期间的一致性。客户端可能按旧路由写，控制面可能已经发布新路由，数据可能正在复制。常见做法是让分片有状态机：准备迁移、双写或转发、追平增量、切换路由、清理旧副本。每一步都要能失败重试。没有状态机的再分片脚本，往往在中途失败后留下难以解释的双写和漏写。

\section{复制}

复制保存多个副本，提高可用性和读吞吐。同步复制延迟高但数据更安全，异步复制延迟低但可能丢失最新写入。主从复制简单，multi-leader 和 leaderless 复制更复杂，需要处理冲突。

复制协议必须说明写入何时算成功。如果 leader 写本地日志就返回，延迟低但 leader 崩溃可能丢数据；如果多数副本确认后返回，延迟高但故障恢复更稳；如果所有副本确认后返回，可用性和尾延迟都受最慢副本影响。工程系统通常让不同数据选择不同级别：元数据强一些，缓存和统计弱一些。

\section{读写仲裁}

Quorum 系统常用 W 个写副本和 R 个读副本满足 R+W>N，从而让读集合和写集合有交集。这个模型帮助理解一致性和可用性权衡。W 大提高写安全性但增加写延迟，R 大提高读新鲜度但增加读延迟。实际系统还要处理慢副本、 hinted handoff、read repair 和版本冲突。

\section{一致性}

强一致性让读者看到最新写入，但通常需要更多同步。最终一致性允许副本暂时不同，换取可用性和低延迟。选择一致性模型要看业务语义：计费、库存和元数据通常更偏强一致；缓存、推荐和日志统计可以接受延迟收敛。

\section{共识}

Raft 和 Paxos 这类共识协议用于在故障中让多个节点就日志顺序达成一致。共识不是让所有操作都变快，而是让关键状态在故障中仍然可恢复。共识路径应尽量小而清晰，不应把所有大数据流量都压进共识日志。

共识适合维护小而关键的顺序状态，例如元数据、配置、leader 选举、分片归属和事务提交点。把大对象、热数据流和高频统计全部放进共识日志，通常会让系统吞吐受限。高性能分布式系统常把控制面和数据面分开：控制面用共识保证关键决策一致，数据面用分片、复制、批处理和校验保证吞吐。

\section{Leader、任期和日志}

以 Raft 的直觉看，系统在一个任期内选出 leader，客户端写入先进入 leader 日志，leader 复制给 follower，多数确认后提交。任期号防止旧 leader 在网络恢复后继续破坏新状态，日志索引和任期共同描述顺序。理解这些概念，能帮助读者读懂很多分布式存储的元数据路径。

\section{读路径和租约}

写路径需要顺序，读路径也要定义新鲜度。强一致读可能需要联系 leader 或多数副本，延迟较高；从 follower 读延迟低，但可能读到旧数据；租约读通过时间约束减少通信，但依赖时钟和租约安全边界。选择读路径时要问：业务是否允许旧读，旧到什么程度，故障切换时如何避免旧 leader 继续服务强一致读。

\section{复制和备份}

复制和备份经常被混为一谈。复制解决的是在线可用性：某个副本不可用时，系统还能从其他副本继续服务。备份解决的是历史恢复：数据被误删、程序 bug 写坏、勒索软件破坏或运维误操作后，系统能回到过去某个正确版本。复制通常追求低延迟和快速同步，备份通常追求隔离、保留周期和可验证恢复。一个系统可以有三副本，仍然必须有备份。

\begin{warningbox}
不要把复制等同于备份。复制会快速传播错误删除和错误写入；备份强调历史恢复。两者解决的问题不同。
\end{warningbox}

\begin{labbox}
实验：设计一个三副本 key-value 存储的写入流程。分别画出同步多数写、异步主从写和读修复流程，并说明各自的延迟和失败行为。
\end{labbox}
